% collatz_ns_conjecture.tex
% Title: The dm³ Operator Grammar: Collatz (TOGT) and Navier–Stokes as Parallel Structural Instances
% G6 LLC · Pablo Nogueira Grossi · Newark NJ · 2026
% Requires: booktabs, amsmath, amssymb, hyperref

\documentclass[11pt,a4paper]{article}

\usepackage{amsmath,amssymb,amsthm}
\usepackage{booktabs}
\usepackage{array}
\usepackage{geometry}
\usepackage{hyperref}
\usepackage{microtype}
\usepackage{lmodern}

\geometry{margin=2.5cm}

\title{The dm\textsuperscript{3} Operator Grammar:\\
       Collatz (TOGT) and Navier--Stokes\\
       as Parallel Structural Instances}
\author{Pablo Nogueira Grossi\\
        G6 LLC, Newark NJ\\
        \texttt{dm3-lab}}
\date{2026}

\theoremstyle{definition}
\newtheorem{defn}{Definition}
\theoremstyle{plain}
\newtheorem{thm}{Theorem}
\newtheorem{conj}{Conjecture}

\begin{document}

\maketitle
\tableofcontents
\clearpage

% ============================================================
\section{Introduction}
% ============================================================

The dm\textsuperscript{3} framework (Discrete/Continuous Manifold Mechanics, level~3) posits a
four-operator generative grammar
\[
  G \;=\; U \circ F \circ K \circ C
\]
that governs both discrete arithmetic dynamical systems and continuous partial-differential
equations. The central claim is that Collatz (via TOGT) and the three-dimensional
Navier--Stokes equations are \emph{instances} of the same grammar, sharing identical operator
roles and a common local-to-global obstruction called \textbf{Bridge~0}.

This document records the formal parallel, states the analytic targets on each side, and
provides supporting LaTeX tables, subsection text, and a Lean~4 formalisation skeleton for
the discrete (Collatz/TOGT) side.

% ============================================================
\section{The dm\textsuperscript{3} Operator Grammar}
% ============================================================

\begin{defn}[dm\textsuperscript{3} generative cycle]
A \emph{dm\textsuperscript{3} generative cycle} is a quadruple of maps
$(C, K, F, U)$ acting on a state space $\mathcal{S}$ such that:
\begin{enumerate}
  \item $C : \mathcal{S} \to \mathcal{S}$ \emph{injects} energy or height (Contact/Compression);
  \item $K : \mathcal{S} \to \mathcal{S} \times \mathbb{R}_{\ge 0}$ \emph{extracts} the
        kernel/valuation, recording local dissipation;
  \item $F : \mathcal{S} \times \mathbb{R}_{\ge 0} \to \mathcal{S}$ \emph{dissipates} flux by
        decrementing the Lyapunov height;
  \item $U : \mathcal{S} \to \mathcal{S}$ \emph{unfolds} back to the next contact point,
        regenerating global structure.
\end{enumerate}
The \emph{macro-step} is $T = U \circ F \circ K \circ C$.
\end{defn}

% ============================================================
\section{TOGT as the Discrete Realisation of dm\textsuperscript{3} Mechanics}
\label{sec:togt-discrete}
% ============================================================

\subsection{TOGT as the Discrete Realisation of dm\textsuperscript{3} Mechanics}

Topographical Orthogonal Generative Theory (TOGT) is the discrete instantiation of the
dm\textsuperscript{3} operator grammar $G = U \circ F \circ K \circ C$. In the Collatz setting
the four operators have precise arithmetic roles: $C$ injects growth via $3n+1$, $K$ extracts
the odd kernel and records the 2-adic valuation, $F$ encodes flux dissipation by decrementing
the discrete Lyapunov height by $v\log 2$, and $U$ folds the result back to the next contact
point. Navier--Stokes supplies the continuous analogue: advection injects energy, scalewise
measurement extracts gradients and vorticity, viscosity dissipates energy, and the return to a
transverse section regenerates the global structure.

The parallel is structural and exact at the level of operator roles and the Bridge~0
obstruction. In both domains a local dissipation mechanism is evident, yet the analytic language
to convert that local control into a global closure theorem is missing. For Navier--Stokes the
missing ingredient is endpoint control of nonlinear fluxes across scales; for TOGT/Collatz it is
a discrete local-to-global dissipation inequality that yields mean contraction. Framing both
problems as instances of the same generative grammar clarifies why progress on the
higher-order closure principle would produce parallel advances across discrete and continuous
settings.

This perspective is empirical and conditional: the polar vortex demonstrates the continuous
dm\textsuperscript{3} structure in nature, and Collatz is the arithmetic test case. The research
program therefore focuses on formalising discrete dm\textsuperscript{3} membership axioms,
calibrating empirical diagnostics, and proving the Bridge~0 closure principle in a higher-order
language that treats continuous and discrete systems as instances of the same operator grammar.

% ============================================================
\section{Side-by-Side Structural Comparison}
\label{sec:table}
% ============================================================

Table~\ref{tab:togt-vs-ns} summarises the exact correspondence between TOGT (discrete
dm\textsuperscript{3}\_disc) and Navier--Stokes (continuous dm\textsuperscript{3}) at each
operator level, together with the shared dm\textsuperscript{3} role and the Bridge~0 parallel.

\begin{table}[ht]
\centering
\small
\begin{tabular}{@{}p{3.2cm}p{4.1cm}p{4.1cm}p{4.1cm}@{}}
\toprule
\textbf{Concept}
  & \textbf{TOGT (Discrete dm\textsuperscript{3}\_disc --- Collatz)}
  & \textbf{Navier--Stokes (Continuous dm\textsuperscript{3})}
  & \textbf{Shared dm\textsuperscript{3} Role / Bridge~0 Parallel} \\
\midrule
Contact Injection
  & $C(n)=3n+1$ (odd-step growth + perturbation)
  & Nonlinear advection $(u\cdot\nabla)u$
  & Local energy / height injection \\[4pt]
Kernel valuation / extraction
  & $K(m)=(m/2^{v_2(m)},\,v_2(m))$ --- records exact dissipation count
  & Extraction of velocity gradient / vorticity at each scale
  & Local dissipation measurement \\[4pt]
Flux dissipation
  & $F$: subtract $v\log 2$ from discrete Lyapunov height
  & Viscous dissipation $\nu\Delta u$ + cascade to small scales
  & Local energy loss / flux control \\[4pt]
Update / unfold
  & $U$: return to next odd contact point; update orbit measure
  & Return to next Poincar\'e section or transverse slice
  & Global structure regeneration \\[4pt]
Macro-step / time evolution
  & $T(n)=U\circ F\circ K\circ C(n)=\dfrac{3n+1}{2^{v_2(3n+1)}}$
  & Evolution under NS PDE over parabolic time scale $\tau_k$
  & One full generative cycle \\[8pt]
Monster threshold / flux barrier
  & $g_6=33$; mean contraction $\mathbb{E}[\log T^2-\log n]<0$ above threshold
  & Shell-flux barrier preventing cascade to small scales
  & Critical stability point \\[4pt]
Bridge~0 obstacle
  & Local 2-adic dissipation + TOGT grammar $\not\Rightarrow$ global mean contraction
  & Local energy inequality + controlled shell flux $\not\Rightarrow$ global regularity
  & The shared local-to-global gap \\[4pt]
Analytic target
  & Mean contraction + Lyapunov descent + structured cycle
  & Flux-barrier $\Rightarrow$ Carleson / KT endpoint $\Rightarrow$ no finite-time singularity
  & Same membership condition in dm\textsuperscript{3} \\
\bottomrule
\end{tabular}
\caption{Side-by-side structural comparison of TOGT (discrete dm\textsuperscript{3}\_disc) and
  Navier--Stokes (continuous dm\textsuperscript{3}). Bridge~0 denotes the local-to-global
  contraction/closure obstacle common to both.}
\label{tab:togt-vs-ns}
\end{table}

% ============================================================
\section{The Discrete Collatz Operators}
\label{sec:operators}
% ============================================================

Let $n \in \mathbb{N}$ be odd. The four TOGT operators act as follows.

\paragraph{Contact injection $C$.}
\[
  C(n) = 3n + 1.
\]
Since $n$ is odd, $C(n)$ is even, initiating the 2-adic cascade.

\paragraph{Kernel extraction $K$.}
Let $v_2(m)$ denote the 2-adic valuation of $m$ (the exponent of the highest power of 2
dividing $m$). Then
\[
  K(m) = \bigl(m/2^{v_2(m)},\; v_2(m)\bigr).
\]
The first component is the odd kernel; the second records how much energy is dissipated in
this step.

\paragraph{Flux dissipation $F$.}
Given a discrete Lyapunov height $h : \mathbb{N}_{\mathrm{odd}} \to \mathbb{R}$, the dissipated
height after one step is
\[
  h(k) \;\longmapsto\; h(k) - v \log 2,\quad (k, v) = K(m).
\]

\paragraph{Unfold $U$.}
$U$ returns the odd kernel $k$ as the next iterate, closing the generative cycle.

\paragraph{Macro-step $T$.}
Composing all four:
\[
  T(n) = U \circ F \circ K \circ C(n) = \frac{3n+1}{2^{v_2(3n+1)}}.
\]
$T$ maps odd integers to odd integers, and the Collatz conjecture asserts that every orbit
under $T$ eventually reaches 1.

% ============================================================
\section{Bridge 0: The Shared Local-to-Global Gap}
\label{sec:bridge0}
% ============================================================

\begin{defn}[Bridge~0 obstacle]
Let $\mathcal{D}$ be a dm\textsuperscript{3} instance (discrete or continuous). A
\emph{Bridge~0 obstacle} is the gap between:
\begin{enumerate}
  \item a provable \emph{local} dissipation or energy inequality at each step of $G$, and
  \item a \emph{global} closure theorem --- mean contraction in the discrete case, global
        regularity in the continuous case.
\end{enumerate}
\end{defn}

\begin{conj}[Bridge~0 closure principle]
There exists a higher-order language $\mathcal{L}$ in which the local dissipation data of any
dm\textsuperscript{3} instance implies the global closure theorem, and in which the Collatz
(TOGT) and Navier--Stokes instances are definitionally equivalent cases.
\end{conj}

The monster threshold $g_6 = 33$ and the mean-contraction estimate
$\mathbb{E}[\log T(n)^2 - \log n] < 0$ for $n > 33$ represent empirical evidence that the
discrete Bridge~0 gap is closeable. The Carleson--Koch--Tataru endpoint regularity theory
represents the current best approach on the continuous side.

% ============================================================
\section{Research Programme}
\label{sec:programme}
% ============================================================

The dm\textsuperscript{3} research programme proceeds along three parallel tracks:

\begin{enumerate}
  \item \textbf{Formalise discrete dm\textsuperscript{3} membership axioms.}
        Define the axiom set for a system to qualify as a dm\textsuperscript{3}\_disc instance,
        including operator roles, Lyapunov height, and the threshold $g_6$.
  \item \textbf{Calibrate empirical diagnostics.}
        Verify the mean-contraction inequality computationally over extended ranges and
        connect it to the formal axioms via the AXLE formalisation (Lean~4).
  \item \textbf{Prove the Bridge~0 closure principle.}
        Develop the higher-order language $\mathcal{L}$ in which local dissipation implies
        global closure, treating discrete and continuous systems as instances of the same grammar.
\end{enumerate}

The polar vortex demonstrates the continuous dm\textsuperscript{3} structure in nature; Collatz
is the arithmetic test case. Progress on either front informs the other through the shared
operator grammar.

% ============================================================
\appendix
\section{Lean 4 AXLE Module: TOGT Operator Skeleton}
\label{app:lean}
% ============================================================

The following Lean~4 pseudocode (module \texttt{TOGT\_Collatz}) provides the AXLE
formalisation skeleton for the discrete TOGT operators. It is part of AXLE Target~5.

\begin{verbatim}
/-
AXLE: TOGT operator skeleton (Lean 4)
This module provides types and basic lemmas for the discrete TOGT operators.
Fill in definitions and proofs as part of AXLE Target 5 formalisation.
-/

structure KernelVal where
  k : Nat   -- odd kernel
  v : Nat   -- 2-adic valuation

def C (n : Nat) : Nat := 3 * n + 1

/-- 2-adic valuation v2(n) -/
def v2 : Nat -> Nat
  | 0 => 0
  | n =>
    let rec aux (m : Nat) (cnt : Nat) : Nat :=
      if m % 2 == 0 then aux (m / 2) (cnt + 1) else cnt
    aux n 0

def K (m : Nat) : KernelVal :=
  { k := m / (2 ^ (v2 m)), v := v2 m }

/-- Discrete Lyapunov height placeholder. -/
def h (k : Nat) : Float := 0.0

/-- Flux dissipation -/
def F (kv : KernelVal) : KernelVal × Float :=
  let k := kv.k
  let v := kv.v
  (kv, h k - Float.ofNat v * Float.log 2.0)

/-- Update / unfold -/
def U (kv_h : KernelVal × Float) : Nat := kv_h.1.k

/-- Macro-step T(n) = U (F (K (C n))) -/
def T (n : Nat) : Nat :=
  let m := C n
  let kv := K m
  U (F kv)

/- Lemmas to prove:
   - T maps odd integers to odd integers (when n is odd)
   - T n = (3n+1) / 2^{v2(3n+1)}
   - Properties of v2 and K
-/

theorem T_eq_formula (n : Nat) :
    T n = (3 * n + 1) / (2 ^ (v2 (3 * n + 1))) := by
  simp [T, U, F, K, C]

/- Bridge 0 schema (axiom template — to be proved, not assumed) -/
-- axiom bridge0_schema :
--   formal statement that local dissipation + TOGT axioms imply mean contraction
\end{verbatim}

The full runnable module (with \texttt{v2} proven and \texttt{T\_eq\_formula} completed) is
maintained in \texttt{lean/TOGT\_Collatz.lean}.

\end{document}
